\section{Finite detection must escape to larger Toeplitz windows}

\begin{theorem}[NN-A014: an explicit finite-window obstruction near the Newman boundary]
Use the normalized coefficients $c_k(t)$ of NN-A007. Let
$T_t^{[R]}=(c_{j-i}(t))_{0\le i,j\le R}$, and for $t<0$ define
\[
 R_*(t)=\min\{R\ge1:T_t^{[R]}\text{ has a negative minor}\}.
\]
Under A this minimum exists for every $t<0$, and
\[
 \lim_{t\uparrow0}R_*(t)=\infty.
\]
No hypothesis on simplicity of zeta zeros is used. The following
explicit quantitative margin is available for every fixed $R$.
Let $\delta_R>0$ be the minimum at $t=0$ of all minors of
$T_0^{[R]}$ not forced to vanish by upper triangularity. Put
\[
 C_R=\max(1,c_1(0),\ldots,c_R(0)),\qquad
 A_R=\max_{0\le k\le R}(2k+2)(2k+1)c_{k+1}(0),
\]
\[
 B_R=(R+1)!(R+1)A_RC_R^R,\qquad
 \epsilon_R=\frac{\delta_R}{2B_R}.
\]
Then all non-forced minors of $T_t^{[R]}$ are at least
$\delta_R/2$ for $-\epsilon_R\le t\le0$, so
$R_*(t)>R$ on $-\epsilon_R\le t<0$.
\end{theorem}
\begin{proof}
NN-A007 gives a negative finite minor at every $t<0$, which is
contained in some initial window. At $t=0$ the positive-zero product
has infinitely many positive parameters. The nonintersecting-index
construction in NN-A008 applies without the normalization $\sum\beta=1$
and proves positivity of every non-forced minor. There are finitely
many minors in a fixed window, so $\delta_R>0$.

The theta integral gives the exact coefficient equation
\[
 c_k'(t)=(2k+2)(2k+1)c_{k+1}(t)-2c_1(t)c_k(t).
\]
Indeed numerator differentiation inserts $u^2$, and denominator
differentiation inserts its normalized mean $2c_1(t)$.
Under the probability measure proportional to
$e^{tu^2}\Phi(u)du$, this derivative is
$\operatorname{Cov}(u^{2k},u^2)/(2k)!$.
For increasing real functions $f,g$, covariance is nonnegative:
with independent $X,Y$ it equals
$\frac12\mathbb E[(f(X)-f(Y))(g(X)-g(Y))]\ge0$.
Thus $0<c_k(t)\le c_k(0)$ for $t\le0$, and
$0\le c_k'(t)\le(2k+2)(2k+1)c_{k+1}(0)$.

An $r$ by $r$ determinant is a sum of $r!$ signed products; its
derivative is a sum of at most $r!r$ products with one differentiated
entry. For $r\le R+1$, each product has absolute value at most
$A_RC_R^{r-1}$, which is bounded by the displayed $B_R$ after
multiplication by $r!r$. Entries below the diagonal and their
derivatives are zero. The fundamental theorem of calculus gives
$|D(t)-D(0)|\le B_R|t|$ for every such minor and $t\le0$.
This proves the stated margin. Since $R$ was arbitrary, it proves
divergence of the minimal detecting window.
\end{proof}

\begin{corollary}[An effective lower bound for the detecting window]
Under A, for every integer $R\ge1$ put
\[
 D_R=\left\lfloor\frac{(R+1)^2}{4}\right\rfloor,
 \qquad
 \underline\epsilon_R=
 \frac{40000^{-D_R}}
 {2^{R+1}(D_R!)^2(R+1)!(R+1)}.
\]
Every minor of $T_t^{[R]}$ not forced to vanish by upper triangularity
is strictly positive for $-\underline\epsilon_R\le t\le0$.
In particular, for all sufficiently small negative $t$,
\[
 R_*(t)\gg
 \sqrt{\frac{\log(1/|t|)}{\log\log(1/|t|)}}.
\]
The implicit constant and the meaning of ``sufficiently small'' are
effective; the displayed factorial bound is an explicit version for
every finite $R$. No simplicity assumption is required.
The later NN-A022 improves the asymptotic lower bound to
$R_*(-a)\gg a^{-1/2}$, with an explicit constant, by controlling
all minors of each bounded order at arbitrary indices.
\end{corollary}
\begin{proof}
First we establish a deliberately generous effective bound for zero
ordinates. Hasanalizade--Shen--Wong, \emph{Counting zeros of the Riemann
zeta function}, arXiv:2107.06506, Corollary~1.2, equation~(1.5), gives
for $T\ge e$
\[
 \left|N(T)-\frac{T}{2\pi}\log\frac{T}{2\pi e}\right|
 \le .1038\log T+.2573\log\log T+9.3675.
\]
The lower bound on the right, after subtracting $T/100$, is positive
at $100$ and strictly increasing for $T\ge100$. For example,
$2\pi<7$, $e<3$, $\log100<5$, and $\log\log100<2$ give its value
at $100$ greater than
$100/7-.104\cdot5-.258\cdot2-9.368-1>0$; its derivative is greater
than $1/7-.104/100-.258/100-.01>0$.
Consequently $N(100j)>j$ and $\gamma_j\le100j$, with ordinates
indexed with multiplicity. Since $x_j=2\gamma_j$ in our heat
normalization, the product parameters satisfy
\[
 b_j=x_j^{-2}\ge\frac{1}{40000j^2}=\frac{\kappa}{j^2}.
\]

Consider an allowed $r$ by $r$ minor with increasing row and column
indices $a_i\le d_i$ in $\{0,\ldots,R\}$. The path construction in
NN-A008 moves the rightmost index first, one unit at a time, and
requires exactly $D=\sum_i(d_i-a_i)$ moves. Here
\[
 D\le r(R+1-r)\le D_R.
\]
Use one move in each of the first $D$ bidiagonal factors
$I+b_jS$, and use diagonal steps in every other factor. Its positive
Cauchy--Binet summand is $\prod_{j=1}^D b_j$; all other summands are
nonnegative. Passing from finite products to the coefficientwise
limit therefore gives
\[
 \delta_R\ge\frac{\kappa^{D_R}}{(D_R!)^2}.
\]
This includes $D=0$, whose chosen summand is $1$.

Uniform coefficient bounds remove the remaining dependence on
unknown finite minima. Define
\[
 C=\sum_{k\ge0}c_k(0)
   =\frac{H_0(i)}{H_0(0)}=\frac{1}{2\xi(1/2)},
 \qquad
 A_0=\sum_{k\ge0}(2k+2)(2k+1)c_{k+1}(0)
   =\frac{\xi''(0)}{4\xi(1/2)}.
\]
The positive theta integral proves convergence of both sums and
$C>1$, $A_0>0$. The accompanying Arb certificate, using
$\xi(s)=(s-1)\pi^{-s/2}\Gamma(1+s/2)\zeta(s)$, rigorously gives
\[
 1<C<2,\qquad 0<A_0<1.
\]
One may retain the exact constants instead of this certificate:
$C=1.00579\ldots$ and $A_0=0.0117395\ldots$.
Since $C_R\le C$ and $A_R\le A_0$, the preceding theorem applies
with derivative bound
$(R+1)!(R+1)2^R$. Substituting the lower bound for $\delta_R$
gives exactly $\underline\epsilon_R$.

For completeness, elementary factorial bounds give, for $R\ge1$,
\[
 \log(1/\underline\epsilon_R)
 \le 10(R+1)^2\log(R+2).
\]
Indeed use $D_R\le(R+1)^2/4$, $\log(n!)\le n\log n$,
$\log40000<11$, and $\log(R+2)>1$ in the explicit expression;
the sum of the resulting coefficients is less than $10$.
Writing $L=\log(1/|t|)$, choose $R+1$ of order
$\sqrt{L/(20\log L)}$. For sufficiently large $L$ the last display
is at most $L/2$, so $|t|<\underline\epsilon_R$ and $R_*(t)>R$.
This proves the asserted effective asymptotic lower bound.
\end{proof}

This is a precise sense in which backward heat leaves the positive
infinite object while every fixed finite truncation initially remains
positive. It does not assert that the minimum \emph{order} of a negative
minor diverges: a bounded-order minor could move to unbounded row or
column indices. The theorem controls the full window parameter and
supplies computable analytic margins, not a polynomial complexity
classification. In representation language, every fixed finite family
of these Toeplitz/Schur tests stays positive for a negative-time interval,
although the full specialization has already left the PF cone. In flag
geometry, every fixed-rank point remains in its positive Toeplitz chart
on that interval; the obstruction to a compatible positive object
occurs at ranks tending to infinity.

\begin{proposition}[A004 refinement: transversality is additional data]
The conditions of a real entire function with all zeros on
$\Re s=1/2$ and functional equation $L(1-s)=L(s)$ do not imply that
the formal shifted-series deformation has a nonzero first normal jet
at every simple zero. An exact example is
\[
 L(s)=((s-1/2)^2+1)(1+\tfrac23(s-1/2)^2),\qquad \rho=1/2+i.
\]
All four zeros are simple and on the line, but
\[
 \rho'(0)=\frac{25i}{4},\qquad
 \rho''(0)=-208+\frac{4979i}{16}.
\]
It therefore moves to the left at second order for small positive
time, despite a purely tangential nonzero first jet.
\end{proposition}
\begin{proof}
In $x=s-1/2$, the zeros are $\pm i,\pm i\sqrt{3/2}$, distinct,
and the polynomial is even with real coefficients.
The series defining $L_t$ converges locally for $|t|<1$, since
$L(s+k)=O(k^4)$ on compact $s$ sets. Substitution into NN-A004 gives
the displayed rational complex numbers; the included symbolic check
uses exact rational arithmetic. Since the first real coefficient is
$\Re(\rho''(0)/2)=-104$, the sign claim follows.
This model has no asserted Euler product. It does not refute
transversality for the actual zeta function; it isolates the extra
phase-nonvanishing input needed for that assertion.
\end{proof}
